\subsection*{重发同一项收款，需要平台认得它}

继续从平台余额 10000、无收款记录开始。购买服务为 C1 选定 \code{provider\_key=pay-C1}。\footnote{\code{pay-C1} 是便于推演的示意值。实际键可以采用足够大的随机标识或不碰撞的受控编号。身份值持久保存，平台访问凭据另由授权配置提供。}\code{provider\_key} 是向平台表达同一次支付操作的键，服务在 A 中保存它，每次发送 C1 都从该记录读取。发送时还要选择平台账户、测试或正式环境及动作。这组解释键的范围保存为 \code{provider\_scope}，示例为 \code{simulator:merchant-7:test:charge.v1}，表示模拟平台的商户 7、测试环境和 \code{charge.v1} 动作。

平台在这个范围内提供幂等受理时，会把 \code{(provider\_scope, provider\_key)} 绑定到第一次有效输入。相同键和输入再次到达，找到原操作；不同输入到达，报告输入冲突。成功记录与实际扣款共同持久保存，才能避免平台内部出现先扣款、后丢失去重记录的缝隙。RIFL 要求完成记录与业务变化原子保存，在这里承担这项责任的是平台自己的提交域~\cite{rifl}。

\CSevenFigure{ch07-original03}{下游身份在平台内部同时关联输入和资金结果}{fig:ch07-provider-identity}

先采用这样一组平台规则：同一有效范围和键至多形成一条成功收款；每次输入按保存版本比较；资金流水、扣后余额与终态响应共同提交；键受保护期间，同参重试取得同一收款编号和原结果。同键并发尝试在平台唯一键处等待，前一个提交就恢复它，前一个回滚则允许后来尝试继续。在这些条件下，\(H_0\) 与 \(H_1\) 都可以重发原 C1：\(H_0\) 第一次完成，\(H_1\) 读取 X1，余额最多减为 7000。

支付输入与购买输入分别保存。购买输入中的 Q7 表达用户接受了哪份报价；发给平台的输入通常包含商户、付款人关联、币种、3000 分、收款动作和本地关联 C1。服务从已经验证的 Q7 取得金额，再按 \code{payment-input-v1} 固定这组支付字段。回调核对平台收到的 3000 与本地快照，重发也使用这份输入。更换付款人、币种或动作，即使整数金额仍为 3000，也改变了收款含义。

同 \(S\) 重试先比较原购买输入，同 C1 重发再保持原支付输入。若恢复时按商品现价构造平台请求，Q7 的 3000 可能变成 Q8 的 3500；有输入比较的平台会报告冲突，缺少比较的平台可能执行被改动的付款。重新报价、改用另一付款工具或重新确认，都要先处理旧 C1 的结局，再按有明确授权的新动作推进。

\subsection*{防止重复和返回什么结果，是两项能力}

幂等受理可以限制效果数量，返回内容则决定调用方能够获知什么。平台可能返回原收款编号和可解释终态，也可能只提示请求已经存在，还可能重放一份没有说明资金结局的错误响应。恢复 P1 时，要逐项判断返回内容证明了什么。

Stripe API v1 提供了一个具体例子。其官方文档规定，在端点执行开始后保存首次状态码和响应体，同键重试返回已保存响应，其中包括 500；执行前的部分校验失败和同键并发冲突采用不同处理~\cite{stripe-v1-idempotency}。因此，同键重试得到相同的 500，带回的是保存的 HTTP 结果。Stripe 的错误处理文档将这类 500 的结局视为不确定，给出通过对象查询和回调核对支付结局的路径~\cite{stripe-v1-idempotency}。

这里引用 API v1 的规则，当前文档另行说明了 API v2 的重试行为。错误是否缓存、保存多久、重试是否继续失败工作，都与具体接口版本有关。前面的平台模型在同键重试时能够恢复稳定成功结果，接入某个端点时，就要把这项假设与它实际返回的内容逐一比较。

结果查询也有定位条件。恢复者持有的身份要能找到操作，返回内容要能关联 C1，并说明它反映哪份平台状态。若查询只接受平台分配的 X1，而创建响应丢失使本地从未取得 X1，这条查询路径此时便用不上。按商户关联号查找、同键重试返回 X1，或事先持久保存可查询的支付对象编号，都可以补上相应的定位路径。

预先创建支付对象，也有响应丢失的位置：平台已经建立对象，购买服务却尚未保存编号。这个创建步骤同样要有稳定身份和恢复办法。平台对象编号一旦可靠保存，就成为后续查询入口；分配该编号时留下的提交缝隙，仍由创建协议处理。

\subsection*{查询在什么条件下代表平台当前事实}

为下面的推演，先采用一种强查询：按完整范围和身份，直接读取平台负责该操作的当前持久记录，覆盖所有受理路径，且已完成记录在保留期内都能被找到。它还满足第六章使用的线性化要求。可以为每次查询选出调用与返回之间的一个读取点，把查询结果解释为那个时点的平台记录。

例如，扣款成功响应已经返回，随后才开始查询，查询就应找到这项仍在保留中的成功。若查询与扣款重叠，读取点可以在 E 之前，准确返回空也符合规则。图~\ref{fig:ch07-query-point} 把两种次序分开画出。

\CSevenFigure{ch07-query-point}{强查询的结果取决于读取点与资金提交的先后}{fig:ch07-query-point}

平台把身份、收款和流水共同提交，查询又读取同一份当前记录时，数据库快照可以充当这个读取点。E 已在快照以前提交，就应读到 X1；E 在快照以后发生，查询可以为空。换成另一家平台，就要重新判断这两个例子是否成立：成功以后何时可查、故障切换后是否保留记录、查询覆盖哪些身份，以及分页和过滤会漏掉什么，都会影响空结果的含义。

实际平台可能同时提供按对象编号读取和搜索。Stripe 的 PaymentIntent 查询用已知编号取得对象详情；其搜索文档明确提示，该搜索不适合要求严格一致性的写后读取，新数据进入搜索结果可能延迟~\cite{stripe-payment-intents}。搜索暂时为空，可能是没有对象，也可能是索引尚未传播。这个提示只对应其搜索接口，不能推成所有平台查询都有相同延迟，也不能据此断言编号读取已经满足本节的全部强查询条件。

异步支付还区分对象存在与资金结局。创建支付对象成功，证明对象已经存在；\code{processing} 表示平台还在处理付款；接口所定义的成功状态才提供付款完成的依据。Stripe 的 PaymentIntent 生命周期区分等待支付方法、需要额外动作、处理中和成功等阶段~\cite{stripe-payment-intents}。本章把一次全额扣款流水提交定义为成功。接入其他支付方式时，适配层要将平台状态与本地 \code{paid} 对应到准确的付款阶段。

\subsection*{两边保留期限怎样接起来}

设本地允许用户在七天内恢复 P1，平台却只在较短窗口内保护 \code{pay-C1}。平台清除键与原响应以后，同字符串重发可能创建第二次收款。购买库继续保留 C1、继续显示等待付款，都没有延长平台的去重期限。

Stripe API v1 文档说明，键至少保留 24 小时，达到期限后可以删除；删除以后再使用同一键，会生成新请求~\cite{stripe-v1-idempotency}。本地据平台规则保存安全发送的时间依据，记为 \code{provider\_retry\_before}。它约束再次发起可能产生效果的调用。超过允许窗口后，旧 P1 的查询可以继续，服务通过查询、回调或对账核查原付款。

\CSevenFigure{ch07-original04}{本地恢复期限和下游安全发送窗口分别约束不同动作}{fig:ch07-windows}

期限的起点也影响判断。平台从首次受理计时，而本地没有获知该时刻时，可以从首次可能发送的本地时刻采用更早的保守截止，并计入时钟误差和已知到达延迟。如果包可以任意延迟，平台又只在收到时判断身份是否受保护，截止以前发出的请求仍可能在清键以后到达。安全发送条件要落在具体机制上：平台永久拒绝过期旧身份、执行端检查固定截止、可证实的在途排空，或已知的到达延迟上限。

把保护截止记为 100。C1 已成功扣款，另一次同键重试在截止前发出，一直滞留途中。平台在 100 以后清理键，101 才收到旧请求。此时本地已经停止新发，平台仍可能再次扣去 3000，使余额从 7000 变为 4000。

若平台支持不可延后的固定截止 100，并在取得执行资格时原子检查它，101 到达的旧请求就会被拒绝，余额保持 7000。图~\ref{fig:ch07-expiry} 对比了两个处理位置。先检查截止、随后另行受理，会给请求留下越过截止的间隔；有效的判期必须与执行资格衔接。

\CSevenFigure{ch07-expiry}{本地停发后，旧请求能否生效取决于执行端的仲裁}{fig:ch07-expiry}

这项拒绝只结束迟到尝试。原 C1 早已成功收取 3000，恢复者继续按原身份查明那项成功，不能把整个 C1 改为拒绝。停止新发送减少新产生的执行机会，查询增加结果证据，已经在途的请求则继续受平台规则约束。

下文对重复发送的论证，都以每个可能生效的请求取得平台执行资格时，固定身份仍受保护为前提。窗口外的未决 C1 保留原映射和输入，转入查询或账单核查。第三章 E1 的关闭控制购买入口，平台对 \code{pay-C1} 的处理由另一套规则决定。
